\epigraph{Análisis del mercado}{%
%<Generalmente este epígrafe se dedica a la sistematización de los fundamentos 
% teórico-metodológicos asociados al tipo de resultado contenido en el campo de acción 
% y que fue reflejado en el objetivo [general] de la investigación; así como establece 
% cuáles de estos fundamentos se constituyeron referentes de la investigación tanto en 
% el contexto nacional como internacional>

    Introducción

    \subepigraph{Panorama General del Mercado}{
        El mercado de sistemas de gestión curricular presenta una segmentación clara
        entre soluciones genéricas de administración académica y herramientas
        especializadas en modelado de variabilidad. Mientras plataformas como
        Moodle, Blackboard y Canvas dominan el segmento de gestión del aprendizaje
        (LMS), existen vacíos significativos en soluciones específicas para la
        planificación curricular basada en Líneas de Productos Educativos (Sánchez et
        al., 2020).
    }

    \subepigraph{Análisis de Soluciones Existentes}{
        \unnumberedsubsection{Sistemas de Gestión Curricular Tradicionales:}
        \begin{itemize}
            \item SIGA y Banner: Los sistemas como Banner son fundamentalmente
            sistemas de información transaccionales diseñados para la gestión
            operativa de una institución. Su arquitectura está optimizada para
            manejar grandes volúmenes de datos de estudiantes, matrículas,
            registros de cursos y pagos. La evidencia de su funcionamiento, como
            se observa en la documentación de "Banner Student", muestra una
            estructura basada en bloques, ventanas y secciones para la gestión de
            admisiones, currículums y campos de estudio (Ellucian Company,
            2015b). Esta lógica está orientada a capturar y recuperar información de
            forma consistente, no a modelar relaciones complejas y variables entre
            componentes curriculares. Carecen de un motor de reglas nativo quepermita definir y verificar restricciones complejas como las que se
            modelan con "requiere" o "excluye" en un Feature Model. Su limitación
            principal es que operan sobre una estructura de datos fija y predefinida,
            lo que los hace rígidos para experimentar con diferentes configuraciones
            curriculares o para gestionar una línea de productos educativos donde
            las relaciones entre asignaturas y competencias son dinámicas y
            variables.
            \item Moodle LMS: es una plataforma excelente para la ejecución y gestión de
            un plan de estudios ya definido. Su fortaleza reside en organizar cursos,
            desglosarlos en unidades y lecciones, gestionar la entrega de
            contenidos, las actividades de evaluación y la interacción entre
            estudiantes y profesores. Sin embargo, su funcionalidad no se extiende
            al diseño instruccional ni a la planificación curricular basada en
            características. Moodle asume que el diseño del curso (los objetivos, la
            secuencia de contenidos, las actividades de evaluación) ya ha sido
            realizado externamente. No proporciona herramientas para que los
            gestores académicos modelen visualmente un plan de estudios, definan
            características opcionales u obligatorias, o establezcan restricciones
            entre ellas. En el contexto de tu investigación, Moodle sería un
            consumidor potencial de la configuración curricular final generada por tu
            backend, pero no una herramienta para crearla.
        \end{itemize}
        \unnumberedsubsection{Herramientas de Modelado de Variabilidad:}
        \begin{itemize}
            \item pure::variants: es una solución comercial robusta para la ingeniería de líneas
            de productos de software (SPL). Proporciona un entorno completo para definir
            características, crear modelos, gestionar restricciones y, crucialmente, derivar
            productos concretos a partir de configuraciones válidas. Su potencia es
            innegable en el dominio técnico para el que fue diseñado. El problema central
            para tu contexto educativo es su alta barrera de entrada para usuarios no
            técnicos. Los gestores académicos y diseñadores curriculares, que son
            expertos en pedagogía pero no necesariamente en SPL, se verían enfrentados
            a una terminología y flujos de trabajo propios de la ingeniería de software.
            Adaptar pure::variants requeriría una personalización significativa de la
            interfaz y una capacitación intensiva, lo que dificulta su adopción ágil en un
            entorno académico, validando la necesidad de una solución más especializaday accesible.
            \item FeatureIDE: es un entorno de código abierto construido como un plugin para
            Eclipse, lo que lo hace muy popular en entornos de investigación y enseñanza
            de la ingeniería de software. Al igual que pure::variants, es poderoso para
            el análisis y la configuración de modelos de características. Sin embargo, su
            dependencia de Eclipse y su orientación hacia desarrolladores lo convierten en
            una herramienta que requiere especialización técnica. Para un coordinador de
            carrera o un vicerrector académico, interactuar con FeatureIDE supondría una
            curva de aprendizaje muy pronunciada. Su arquitectura no está diseñada para
            integrarse de manera natural con los sistemas de información universitarios
            existentes (como los propios SIGA o Banner), lo que limita su utilidad práctica
            para la gestión curricular del día a día. Su uso en educación se limita
            predominantemente a la enseñanza de conceptos de SPL, no a la gestión
            operativa de currículos.
        \end{itemize}
    }

    \subepigraph{Oportunidades de Mercado Identificadas}{
        El mercado asociado a la temática padece de un vacío crítico. No existen
        soluciones que integren efectivamente la gestión de variabilidad de SPL con la
        planificación curricular educativa, creando una oportunidad para sistemas
        especializados (Sánchez et al., 2020). Las herramientas existentes de SPL están
        diseñadas para ingenieros de software, no para profesionales de la educación,
        limitando su adopción en instituciones educativas (Chen et al., 2020).

        La solución propuesta se posiciona estratégicamente en el mercado educativo
        como un "Motor de Variabilidad Curricular" especializado, que capitaliza las
        tendencias actuales de crecimiento en educación híbrida y multimodal, así
        como las demandas crecientes de personalización educativa y eficiencia en la
        gestión académica post-pandemia. Su propuesta de valor se sustenta en
        fortalezas técnicas diferenciadoras, incluyendo un backend especializado en
        gestión de variabilidad curricular, un motor de reglas educativo que comprende
        prerrequisitos y competencias, una API RESTful para integración con sistemas
        existentes y un modelo de datos optimizado para planificación académica.
        Estas capacidades se traducen en ventajas competitivas tangibles como la
        automatización de validaciones que reduce errores en diseño curricular, la
        generación de múltiples variantes de planes de estudio desde un tronco comúny su 
        adaptabilidad a contextos institucionales diversos, posicionándose como
        complemento esencial -no competencia- de los sistemas LMS existentes, al
        enfocarse específicamente en el diseño y planificación académica más que en
        la ejecución curricular.
    }

}
